% =============================================================================
%  s05-wrapup.tex —— 原文 §5 Tying Up Loose Ends（p.32–36）与 References（p.36–38）
%  编号公式：(5.1)–(5.4)
% =============================================================================

% ==== 原文 p.32 ====
\section{收束：未竟之结}\label{sec:5}
% 原文标题：Tying Up Loose Ends

% [Tutor Note: 本节重点] 物理图像 / 数学工具 / 主要结论
\begin{keytakeaways}
\begin{itemize}[leftmargin=1.4em,itemsep=2pt,topsep=2pt]
  \item \textbf{物理图像：}三维体里的陈--西蒙斯理论与它边界上的二维共形场论是同一套数据的两种说法：体中的威尔逊线在边界上化成原初场的算符插入，这正是全息对偶的原型。把规范群换成非紧的庞加莱群，三维引力本身也变成一个陈--西蒙斯理论。
  \item \textbf{数学工具：}琼斯多项式及其 skein 关系与变量 $q$、连通和下配分函数的乘法性质 \eqref{eq:4.7}、Segal 的模函子 \cite{segal} 与共形块、韦斯--朱米诺--维滕模型的作用量、二次卡西米尔 $c_2(G)$。
  \item \textbf{主要结论：}$G=\mathrm{SU}(N)$ 时正确变量是 $q=\exp[2\pi i/(k+N)]$（只有树图级才是 $e^{2\pi i/k}$）；无结的威尔逊圈期望值为 $(q^{N/2}-q^{-N/2})/(q^{1/2}-q^{-1/2})$；一圈修正等价于 $k\to k+c_2(G)/2$，对应二维 CFT 的中心荷 $c=kd/(k+N)$。
\end{itemize}
\end{keytakeaways}

\paragraph{统一的视角}
% 原文标题：A unified view.
在这篇综述里，我们试着理清陈--西蒙斯理论背后潜伏的那些拓扑与几何。旅程从 \S\ref{sec:2} 开始：我们先请出 $\mathrm{U}(1))$ 理论，视角或许取自一位初次撞见其怪异物理的实验物理学家。这段穿越阿贝尔理论的现象学漫游，为随后用路径积分（\S\ref{sec:3}）和正则量子化（\S\ref{sec:4}）对非阿贝尔理论作更形式化的处理做了铺垫。一路上，我们刻意不谈 CS 理论所处的那个更大的物理与数学世界中的许多侧面。下面我们先对已经走过的这条窄路作最后一次回顾，再匆匆一瞥那个更大世界的若干角落，以此作结。

\subsection{全景回顾}\label{sec:5.1}
% 原文标题：Broad Recapitulation

\paragraph{阿贝尔理论}
在 \S\ref{sec:2} 中我们发现，CS 理论是拓扑的。它最简单的物理可观测对象是任意子交换相位，这些相位度量的是附着在平面上游走电荷上的磁通线绕了多少圈。我们很快认出，这些相位就是威尔逊圈的期望值，而用路径积分处理威尔逊圈最为自然。纯 CS 路径积分告诉我们：有磁通存在时，为保住量子理论的规范不变性，陈--西蒙斯能级 $k$ 必须量子化成整数。
% ==== 原文 p.33 ====
带源的 CS 路径积分则表明，任意子相位真正抓住的是拓扑不变量——源所附威尔逊圈扫出的那些纽结的环绕数。此后我们用正则形式处理 $\mathrm{U}(1))$ 理论，发现哈密顿量恒为零，CS 规范场的各个分量彼此正则共轭。我们在 $\R^2$ 上精确解出了这个理论，只得到一个物理态；讨论了它在 $T^2$ 上有限维的可观测代数；最后以若干有趣而重要的拓扑评述收尾。

\paragraph{路径积分}
搭起非阿贝尔规范理论的机器之后，我们引入并量子化了非阿贝尔 CS 理论，入手处是考察规范不变性在绕数非零的规范变换下如何失效。接着我们尝试在弱耦合下计算纯 CS 路径积分。这一计算给出底流形 $M$ 的一大批拓扑不变量，可结果看上去仍依赖它的度规。我们设法消掉了这份依赖，代价是加一个依赖 $M$ 标架的抵消项。最终我们得到纯 CS 理论配分函数的一个显式拓扑公式，还有一条描述它在标架变换下如何变化的规律。带威尔逊圈的路径积分有着相似的拓扑命运：为消去与威尔逊圈自环绕数相伴的发散，我们必须给纽结选定一个标架。与无源理论的类比十分紧密，我们得到一个公式，描述威尔逊圈期望值在标架变换下的变换。

\paragraph{正则量子化}
要对非阿贝尔 CS 理论作正则分析，我们把一般的三维流形 $M$ 切成若干局部形如 $\Sigma\times\R$ 的块，其中 $\Sigma$ 是黎曼面。CS 理论在 $\Sigma$ 上的经典相空间是 $\Sigma$ 上平坦联络的模空间；为量子化它，我们发展出凯勒量子化这一有力技术。这一技术把纯 CS 希尔伯特空间描述成 $M$ 上某个线丛的全纯截面空间。我们还考虑了威尔逊圈的插入：这些圈穿透 $\Sigma$，在其表面留下非阿贝尔电荷，位置是特殊点（distinguished points），也叫标记点，与阿贝尔理论里的磁通线遥相呼应。这些圈起初似乎带来无法解决的问题，连理论的经典描述也躲不开。不过，借助博雷尔--韦伊--博特定理，再靠「带对称性的系统如何量子化」这一哲学，我们微调模空间，把非阿贝尔电荷的效应算了进来。在 $\Sigma$ 亏格为零这一具体情形，能拿到的结果要具体得多，包括希尔伯特空间的完整描述及其维数。

\subsection{延伸与联系}\label{sec:5.2}
% 原文标题：Extensions and Connections

\paragraph{琼斯多项式}
Witten 那篇奠基性论文 \cite{witten1989} 的主要贡献之一，是把一个重要的纽结不变量——琼斯多项式——用 CS 理论里的威尔逊圈来解释。为陈述他的部分结果，我们先看一个在全文里无处不在的量：$e^{2\pi i/k}$。我们给它取名 $q$，并断言：从纽结理论的角度看，这正是写下最终结果所该用的变量。
% ==== 原文 p.34 ====
对复的 $k$，变量 $q=e^{2\pi i/k}$ 与椭圆函数论中出现的 nome $q=e^{2\pi i\tau}$ 关系密切。

椭圆函数与琼斯多项式 $V_q$ 都在 $k=1/\tau$（即 CS 耦合的倒数）取整数值时表现出特殊性质，这绝非巧合：量子 CS 理论的结果恰在此时变得可用。其实，这话只在所谓「树图级」上才成立：对 $G=\mathrm{SU}(N)$，正确的变量其实是 $q=\exp\left(\frac{2\pi i}{k+N}\right)$。用 nome 表示，我们终于能写下非阿贝尔 CS 理论中若干威尔逊圈的期望值。例如 $S^3$ 中一个无结的威尔逊圈 $W[\bigcirc]$，取 $\mathrm{SU}(N)$ 的基础表示，其真空期望值为
\begin{equation}\tag{5.1}\label{eq:5.1}
  \langle W[\bigcirc]\rangle
  = \frac{q^{N/2}-q^{-N/2}}{q^{1/2}-q^{-1/2}}
  \overset{N=2}{\sim} q^{1/2}+q^{-1/2}
  = \frac{\sin[2\pi/(k+2)]}{\sin[\pi/(k+2)]} .
\end{equation}

% [Tutor Note: 易错点] (5.1) 的幂次：分子为 $q^{N/2}-q^{-N/2}$、分母为 $q^{1/2}-q^{-1/2}$，且 $q$ 必须取 $q=\exp[2\pi i/(k+N)]$；若用树图的 $q=e^{2\pi i/k}$，则 $q^{1/2}+q^{-1/2}=2\cos(\pi/k)$ 与右端的 $2\cos[\pi/(k+2)]$ 不等。
\begin{yizhu}[易错点]
\eqref{eq:5.1} 的幂次要逐个核对：分子是 $q^{N/2}-q^{-N/2}$，分母是 $q^{1/2}-q^{-1/2}$，指数上的 $1/2$ 一个都不能漏。更要紧的是 $q$ 的取法：右边化成 $\sin$ 比值时用的是 $q=\exp[2\pi i/(k+N)]$；在 $N=2$ 这一支上，$q^{1/2}+q^{-1/2}=2\cos[\pi/(k+N)]$，恰等于 $\sin[2\pi/(k+N)]/\sin[\pi/(k+N)]$。若照前文的树图定义 $q=e^{2\pi i/k}$ 代入，得到 $2\cos(\pi/k)$，与 $\sin[2\pi/(k+2)]/\sin[\pi/(k+2)]=2\cos[\pi/(k+2)]$ 不等。作者在紧邻 (5.1) 之前的那句话里承认这「只在树图级上」才对，正指此处。
\end{yizhu}

% [Tutor Note: 物理图像] (5.1) 的比值是量子群 $U_q(\mathfrak{sl}_N)$ 中基础表示的量子维数 $[N]_q$，$k\to\infty$ 时回到经典维数 $N$。
\begin{yizhu}[物理图像]
\eqref{eq:5.1} 的比值不是拼凑出来的：它就是量子群 $U_q(\mathfrak{sl}_N)$ 中基础表示的量子维数 $[N]_q$，也就是无结威尔逊圈在 CS 理论里能取到的「大小」。$k\to\infty$（即 $q\to1$）时 $[N]_q\to N$，与经典维数吻合；$k$ 有限时它才明显偏离 $N$，有限能级的量子效应第一次以纯数值的形式露了面。
\end{yizhu}

代入著名的 skein 关系（Witten 也用希尔伯特空间方法导出了它），这个表达式给出无结的正确的（平凡的！）琼斯多项式：
\begin{equation}\tag{5.2}\label{eq:5.2}
  -q\,V_q(\bigcirc) + \left(q^{1/2}-q^{-1/2}\right)\langle W\rangle + q^{-1}V_q(\bigcirc) = 0
  \quad\Longrightarrow\quad V_q(\bigcirc) = 1 .
\end{equation}

% [Tutor Note: 易错点] (5.2) 三个圈同形，原文未标 $L_\pm$ 与 smoothing：标准形为 $q^{-1}V(L_+)-qV(L_-)=(q^{1/2}-q^{-1/2})V(L_0)$，故首项应读作 $L_-$、末项读作 $L_+$、中间项读作 smoothing；且 (5.2) 的整体符号与标准形相反。
\begin{yizhu}[易错点]
\eqref{eq:5.2} 里三个圈都写作 $\bigcirc$，原文没有标注哪两个是 $L_\pm$、哪一个是 smoothing。标准的 Jones skein 关系是 $q^{-1}V(L_+)-qV(L_-)=(q^{1/2}-q^{-1/2})V(L_0)$；据此，首项 $-qV_q(\bigcirc)$ 应读作 $L_-$，末项 $q^{-1}V_q(\bigcirc)$ 应读作 $L_+$，中间的 $\langle W\rangle$ 才是 smoothing 后得到的链环。还要留意符号：标准归一化下二分量平凡链环的 $V(L_0)=-(q^{1/2}+q^{-1/2})$，而 \eqref{eq:5.1} 在 $N=2$ 时取正号，所以 \eqref{eq:5.2} 隐含的符号约定与标准形差一个整体符号。核对时以「把 \eqref{eq:5.1} 代进去后化为恒等式」为准，不要逐项对号。
\end{yizhu}

事实上，琼斯多项式的若干初等性质，从 \S\ref{sec:4.4} 的讨论里已经能看出来。比如，因子化结果 \eqref{eq:4.7} 可以写成
\begin{equation}\tag{5.3}\label{eq:5.3}
  \frac{\mathcal{Z}(M)}{\mathcal{Z}(S^3)}
  = \frac{\mathcal{Z}(M_1)}{\mathcal{Z}(S^3)}\cdot\frac{\mathcal{Z}(M_2)}{\mathcal{Z}(S^3)} .
\end{equation}
把威尔逊圈也算进来时，上面出现的那些比率就变成琼斯多项式；公式于是表达了它们在连通和下的乘性行为。

% [Tutor Note: 易错点] (5.3) 的乘性只对分处连通和两片的链环成立；比率用 $\mathcal{Z}(S^3)$ 归一化正是为了与 $V_q(\text{无结})=1$ 对齐，而同在 $S^3$ 里不相交的两个平凡圈各给出因子 $-(q^{1/2}+q^{-1/2})$，不是 $1$。
\begin{yizhu}[易错点]
\eqref{eq:5.3} 的乘性不可越界使用：它说的是连通和 $M_1\#M_2$ 的配分函数等于两片各自配分函数之积再除以一个 $\mathcal{Z}(S^3)$，对应的链环分处两片。若把两个圈放进同一个 $S^3$，Jones 多项式要对每个分裂出的平凡分量各乘一个因子 $-(q^{1/2}+q^{-1/2})$，结果不是 $1$。比率用 $\mathcal{Z}(S^3)$ 归一化，恰好让空流形对应 $1$、无结对应 $V_q(\bigcirc)=1$；把这两件事混为一谈，是本节最容易犯的错。
\end{yizhu}

\paragraph{与二维 CFT 的关系}
% 原文标题：Relation to 2D CFT.
到此为止，我们一直刻意不谈 CS 理论与二维共形场论之间的深刻联系。完整处理这些联系超出本综述的范围，但不妨在此收集几条提示性的评论。第一条线索来自 \S\ref{sec:3.2} 中我们对纯 CS 理论配分函数的路径积分计算。领头阶结果 \eqref{eq:3.10} 与一圈结果 \eqref{eq:3.27} 的主要差别，除复杂度增加之外，就是组成配分函数的那些相对相位整体平移了，仿佛我们已经把 $k$ 换成 $k+\frac{c_2(G)}{2}$：
\begin{equation}\tag{5.4}\label{eq:5.4}
  \exp\left(ik\,I[A^{(\alpha)}]\right)
  \;\leadsto\;
  \exp\left[i\left(k+\frac{c_2(G)}{2}\right)I[A^{(\alpha)}]\right] .
\end{equation}

% [Tutor Note: 易错点] (5.4) 的 $k\to k+c_2(G)/2$ 与紧随其后的 $c=kd/(k+N)$ 靠 $c_2(G)$ 的归一化配平：作者用 $c_2(\mathrm{SU}(2))=4$，故 $c_2(G)/2=h^\vee$，$\mathrm{SU}(N)$ 时即 $N$；换用 $\Tr_{\rm fund}(T^aT^b)=\frac12\delta^{ab}$ 的归一化则平移应写成 $k+c_2(G)$。
\begin{yizhu}[易错点]
\eqref{eq:5.4} 的平移 $k\to k+c_2(G)/2$ 与紧随其后的中心荷公式 $c=\frac{kd}{k+N}$ 看似差一个因子，其实靠 $c_2(G)$ 的归一化配平：作者在 \eqref{eq:3.32} 之后注明 $c_2(\mathrm{SU}(2))=4$，可见他把 $c_2(G)$ 归一化成 $2h^\vee$（$h^\vee$ 为对偶考克斯特数），于是 $c_2(G)/2=h^\vee$，$\mathrm{SU}(N)$ 时恰为 $N$。这正好是 $q=\exp[2\pi i/(k+N)]$ 与 $c=\frac{kd}{k+N}$ 里的那个 $N$。若改用物理文献常见的 $\Tr_{\rm fund}(T^aT^b)=\frac12\delta^{ab}$ 归一化，$c_2(\mathrm{SU}(2))=2$，平移要写成 $k+c_2(G)$；照字面读成 $k+1$ 会差一倍。
\end{yizhu}

这让人想起二维 CFT 中的类似现象：只要把近似公式里的参数平移，许多半经典结果就变成精确的。实际上，当 $G=\mathrm{SU}(N)$ 时，$k$ 与某个二维 CFT 的中心荷 $c$ 由公式 $c=\frac{kd}{k+N}$ 相联系，其中 $d=\dim G=N^2-1$。CS 耦合弱时 $k\to\infty$，于是 $c\to d$。这样一来（从 CFT 的观点看）毫不奇怪：
% ==== 原文 p.35 ====
CS 理论里在大 $k$ 下导出的许多结果，例如配分函数的标架变换规律 \eqref{eq:3.28}，都有把 $d$ 换成 $c$ 得到的推广。威尔逊圈的标架变换规律 \eqref{eq:3.31} 上演着相似的故事；这一回，因子 $\frac{n_a^2}{2k}$ 换成了某个 CFT 原初场的共形权重 $h$。更一般地，CS 理论中标架化这一概念，在相应 CFT 的单值性里有直接的对应物。

\paragraph{全息对偶}
% 原文标题：Holographic duality.
读到这里应该清楚：三维的 CS 理论与二维的共形场论关系密切。我们握有大量证据：其一，上面那些类比本身就足够撩人。其二，在 \S\ref{sec:4} 中我们瞥见：CS 理论的希尔伯特空间与配分函数，来自沿黎曼面切分三维流形，并仔细照料这些低维边界上的物理。我们在那里遇到的构造——黎曼面的模空间、平坦联络的模空间、全纯向量丛、单值性、表示的融合——都是二维 CFT 的重要配料。其中或许最重要的，是 Segal 的模函子 \cite{segal}。在 Segal 的原始表述里，一张黎曼面 $\Sigma$ 定义了二维 CFT 所栖身的时空；与这个时空对应的是一个向量空间 $\mathcal{H}_\Sigma$，由单位算符的后代所满足的共形沃德恒等式的解组成。Segal 对这个「共形块空间」的构造，与我们对 $\Sigma\times\R$ 上 CS 希尔伯特空间的构造完全同一；这就把 CS 与 CFT 两种理论的等价牢牢确立下来。确实可以验证：我们算出的带源 CS 希尔伯特空间的维数，与对应带标记点黎曼面上二维 CFT 的共形块数目一致。

% [Tutor Note: 含糊点] 「完全同一」「牢牢确立」跳过了要验证的东西：Segal 模函子 [24] 与 CS 几何量子化给出同一组共形块，须核对维数（Verlinde [28]、Moore--Seiberg [29,30]）、拼接相容性与单值性相位三件事。
\begin{yizhu}[含糊点]
「完全同一」与「牢牢确立」两句跨过了实质验证。模函子（Segal \cite{segal}）是从带标记黎曼面到向量空间的函子，配以拼接相容性：它把 $\Sigma$ 映成共形块空间，把拼接映成线性映射的复合。CS 一侧给出的则是平坦联络模空间上的全纯截面（几何量子化）。要说两者等同，至少要核对三件事：维数公式一致（融合规则与 Verlinde 公式 \cite{verlinde1988}，Moore--Seiberg \cite{moore1989,moore1988} 的讨论）、拼接映射与 CS 粘合相容、标架与单值性的相位一致。原文只点出了亏格零情形的维数吻合。另可留意：它把 $\mathcal{H}_\Sigma$ 说成「单位算符的后代」的解空间，这其实是真空（无插入）情形的说法；带标记点时须换成相应原初场的后代，紧接着的维数断言才有着落。
\end{yizhu}

现在轮廓渐渐显出的图景是一张全息图：三维流形 $M$ 上的 CS 理论，等价于（或说对偶于）栖身在黎曼面 $\Sigma$ 上的某个二维 CFT，而 $\Sigma$ 可以看作 $M$ 的「边界」。（在我们关于拼接与手术的讨论里，这个解释是我们硬按上去的。）相应的 CFT 叫韦斯--朱米诺--维滕（WZW）模型，它的作用量与 CS 作用量变分中那个边界项 \eqref{eq:3.4} 颇为相像。在这个框架里，CS 威尔逊线穿透 $\Sigma$ 处的标记点，代表 WZW 模型里原初场的算符插入，这些原初场按威尔逊线所附的表示变换。CS--WZW 对应是一条真正的定理，可以看作全息原理的原型。

% [Tutor Note: 含糊点] 「颇为相像」可落到实处：(3.4) 的第二项 $\frac{k}{4\pi}\int\frac13\Tr[(g^{-1}\partial g)^3]$ 正是 WZW 作用量的 Wess--Zumino 项，即纯规范 $A=g^{-1}\dd g$ 的 CS 作用量在 3 维球上的取值。
\begin{yizhu}[含糊点]
「颇为相像」不是修辞，落到实处是这样：\eqref{eq:3.4} 的第二项 $\frac{k}{4\pi}\int_M\frac{1}{3}\Tr[(g^{-1}\partial g)^3]$ 正是 WZW 作用量的 Wess--Zumino 项。把二维边界 $\Sigma$ 换成以它为边界的 3 维球 $B$，并把 $g$ 延拓进 $B$，该项就是 $\frac{k}{12\pi}\int_B\Tr(g^{-1}\dd g)^3$，即纯规范 $A=g^{-1}\dd g$ 的 CS 作用量。换言之，WZW 作用量 = $\Sigma$ 上的动能项 + CS 作用量在纯规范上的取值；CS--WZW 对应在作用量层面就体现在这里。\eqref{eq:3.4} 的第一项是全导数，搬到边界上化为 $\Sigma$ 的边界积分。
\end{yizhu}

\paragraph{三维引力}
% 原文标题：3D gravity.
有一个重要因素使 CS--WZW 算不上全息对偶的真实例子：对应关系的 CS 一侧没有引力物理。然而 Witten 表明 \cite{witten1988}，三维引力就是一个陈--西蒙斯理论！广义相对论的爱因斯坦场方程在 $2+1$ 维不允许引力波解。因此这个理论没有局部传播自由度；正因如此，它是纯拓扑的。支配该理论的爱因斯坦--希尔伯特作用量有标架场与自旋联络的一阶表述，Witten 表明它可以改写成形如 \eqref{eq:3.3} 的 CS 作用量。三维引力里的「规范场」取值于庞加莱群 $\mathrm{ISO}(2,1)$ 的李代数，或它的某个非紧变体；
% ==== 原文 p.36 ====
具体取哪个，取决于宇宙学常数的符号。（对德西特背景，群是 $\mathrm{SO}(3,1)$；对反德西特，是 $\mathrm{SO}(2,2)$。）由于规范群不再紧致，紧致李群理论的许多内容不再适用，理论也就更难解释。尽管有这份困难，Witten 仍论证 \cite{witten2007}，三维引力容许一个良定义的量子化，与我们在本文所见 CS 理论的量子化类似。这一结论把 CS--WZW 对应确立为 $\mathrm{AdS}_3/\mathrm{CFT}_2$ 对偶的一个真实化身。Witten 还论证该理论无论经典还是量子都精确可解。即便如此，关于这个对偶及其对量子引力的含义，仍有太多东西有待理解。

% [Tutor Note: 含糊点] 「无引力波 \$\\ 无局部自由度 \$\\ 纯拓扑」跳过了一步：2+1 维度规 6 个分量减 3 个微分同胚、3 个约束得 0 个局部自由度；拓扑性来自 Witten [31] 的一阶表述，标架场与自旋联络并成 ISO(2,1)（或 SO(3,1)/SO(2,2)）值联络，作用量即 CS 形式且不含度规。
\begin{yizhu}[含糊点]
「没有引力波 \$\\ 没有局部传播自由度 \$\\ 纯拓扑」跳过了中间一步。2+1 维度规有 6 个分量，减去 3 个微分同胚自由度与 3 个（1 个哈密顿 + 2 个动量）约束，局部自由度恰好为 0；这才是计数上的理由。拓扑性的来源则是 Witten \cite{witten1988} 的一阶表述：标架场 $e^a$ 与自旋联络 $\omega^a$ 并成 $\mathrm{ISO}(2,1)$（或 $\mathrm{SO}(3,1)$、$\mathrm{SO}(2,2)$）值联络 $\mathcal{A}=e^aP_a+\omega^aJ_a$，爱因斯坦--希尔伯特作用量 $\int e\wedge(d\omega+\omega\wedge\omega)$ 就是 \eqref{eq:3.3} 形状的 CS 作用量；作用量里不含度规，故不依赖度规。非紧群使 \cite{witten1989} 那套紧致李群论证失效，这正是 \cite{witten2007} 要处理的技术难点。
\end{yizhu}

% ==== 原文 p.36 下半页：References ====
\begin{thebibliography}{32}

\bibitem{witten1989} E. Witten, ``Quantum Field Theory and the Jones Polynomial,'' \emph{Commun. Math. Phys.} \textbf{121}, 351--399 (1989) doi:10.1007/BF01217730

\bibitem{dunne} G. V. Dunne, ``Aspects of Chern-Simons theory,'' [arXiv:hep-th/9902115 [hep-th]]

\bibitem{tong_gauge} D. Tong, ``Gauge Theory,'' ch. 8, damtp.cam.ac.uk/user/tong/gaugetheory/gt.pdf

\bibitem{tong_qhe} D. Tong, ``Lectures on the Quantum Hall Effect,'' ch. 5, [arXiv:1606.06687 [hep-th]]

\bibitem{chern1974} S.-S. Chern and J. Simons, ``Characteristic forms and geometric invariants,'' \emph{Ann. Math.} \textbf{99} (1), 48--69 (1974) doi:10.2307/1971013

\bibitem{polyakov1988} A. M. Polyakov, ``Fermi--Bose Transmutations Induced by Gauge Fields,'' \emph{Mod. Phys. Lett. A} \textbf{3}, 325 (1988) doi:10.1142/S0217732388000398

\bibitem{ricca2011} R. L. Ricca and B. Nipoti, ``Gauss' Linking Number Revisited,'' \emph{J. Knot Theory Ramif.}, \textbf{20} (10), 1325--1343 (2011) doi:10.1142/S0218216511009261

\bibitem{achucarro1986} A. Achucarro and P. K. Townsend, ``A Chern-Simons Action for Three-Dimensional anti-De Sitter Supergravity Theories,'' \emph{Phys. Lett. B} \textbf{180}, 89 (1986) doi:10.1016/0370-2693(86)90140-1

\bibitem{steenrod1951} N. Steenrod. \emph{The Topology of Fibre Bundles (PMS-14)}, Princeton University Press (1951) JSTOR, \href{https://www.jstor.org/stable/j.ctt1bpm9t5}{www.jstor.org/stable/j.ctt1bpm9t5}

\bibitem{nakahara} M. Nakahara, ``Geometry, Topology and Physics,'' ch. 9--10

\bibitem{coleman} S. Coleman, ``Aspects of Symmetry: Selected Erice Lectures,'' \S7.3.2 doi:10.1017/CBO9780511565045

\bibitem{bott1956} R. Bott, ``An Application of the Morse Theory to the Topology of Lie Groups,'' \emph{Bulletin de la S.M.F.} \textbf{84}, 251--281 (1956), doi:10.24033/bsmf.1472, \href{http://www.numdam.org/articles/10.24033/bsmf.1472}{http://www.numdam.org/articles/10.24033/bsmf.1472}

% ==== 原文 p.37 ====
\bibitem{barnatan1991} D. Bar-Natan and E. Witten, ``Perturbative expansion of Chern-Simons theory with noncompact gauge group,'' \emph{Commun. Math. Phys.} \textbf{141}, 423--440 (1991) doi:10.1007/BF02101513

\bibitem{schwarz1978} A. S. Schwarz, ``The Partition Function of Degenerate Quadratic Functional and Ray-Singer Invariants,'' \emph{Lett. Math. Phys.} \textbf{2}, 247--252 (1978) doi:10.1007/BF00406412

\bibitem{aps1975} M. F. Atiyah, V. K. Patodi, and I. M. Singer, ``Spectral asymmetry and Riemannian Geometry 1,'' \emph{Math. Proc. Cambridge Phil. Soc.} \textbf{77}, 43 (1975) doi:10.1017/S0305004100049410

\bibitem{aps1976a} M. F. Atiyah, V. K. Patodi, and I. M. Singer, ``Spectral asymmetry and Riemannian geometry 2,'' \emph{Math. Proc. Cambridge Phil. Soc.} \textbf{78}, 405 (1976) doi:10.1017/S0305004100051872

\bibitem{aps1976b} M. F. Atiyah, V. K. Patodi, and I. M. Singer, ``Spectral asymmetry and Riemannian geometry 3,'' \emph{Math. Proc. Cambridge Phil. Soc.} \textbf{79}, 71--99 (1976) doi:10.1017/S0305004100052105

\bibitem{parker1984} P. E. Parker, ``On Some Theorems of Geroch and Stiefel,'' \emph{J. Math. Phys.} \textbf{25}, 597 (1984) doi:10.1063/1.526209

\bibitem{milnor1974} J. W. Milnor and J. D. Stasheff, \emph{Characteristic Classes} (AM-76) (1974) doi:10.1515/9781400881826

\bibitem{elhamdadi2020} M. Elhamdadi, M. Hajij, and K. Istvan, ``Framed Knots,'' \emph{Math. Intell.} \textbf{42}, 7--22 (2020) \href{https://link.springer.com/content/pdf/10.1007/s00283-020-09990-0.pdf}{https://link.springer.com/content/pdf/10.1007/s00283-020-09990-0.pdf}

\bibitem{atiyahbott1982} M. F. Atiyah and R. Bott, ``The Yang-Mills equations over Riemann surfaces,'' \emph{Phil. Trans. Roy. Soc. Lond. A} \textbf{308}, 523--615 (1982)

\bibitem{woodhouse} N. Woodhouse, ``Geometric Quantization.''

\bibitem{quillen1986} D. Quillen, ``Determinants of Cauchy--Riemann operators over a Riemann surface,'' \emph{Fuct. Anal. Appl.} \textbf{19}, 31 (1986)

\bibitem{segal} G. Segal, ``Conformal field theory,'' Oxford preprint; and lecture at the IAMP Congress, Swansea, July, 1988

\bibitem{connes2006} A. Connes and M. Marcolli, ``A Walk in the noncommutative garden,'' [arXiv:math/0601054 [math.QA]].

\bibitem{kirillov2004} A. A. Kirillov, \emph{Lectures on the Orbit Method}, American Mathematical Soc., \textbf{64} (2004) doi:10.1090/gsm/064

% ==== 原文 p.38 ====
\bibitem{beasley2013} C. Beasley, ``Localization for Wilson Loops in Chern-Simons Theory,'' \emph{Adv. Theor. Math. Phys.} \textbf{17}, no.1, 1-240 (2013) doi:10.4310/ATMP.2013.v17.n1.a1 [arXiv:0911.2687 [hep-th]].

\bibitem{verlinde1988} E. P. Verlinde, ``Fusion Rules and Modular Transformations in 2D Conformal Field Theory,'' \emph{Nucl. Phys. B} \textbf{300}, 360-376 (1988) doi:10.1016/0550-3213(88)90603-7

\bibitem{moore1989} G. W. Moore and N. Seiberg, ``Classical and Quantum Conformal Field Theory,'' \emph{Commun. Math. Phys.} \textbf{123}, 177 (1989) doi:10.1007/BF01238857

\bibitem{moore1988} G. W. Moore and N. Seiberg, ``Polynomial Equations for Rational Conformal Field Theories,'' \emph{Phys. Lett. B} \textbf{212}, 451-460 (1988) doi:10.1016/0370-2693(88)91796-0

\bibitem{witten1988} E. Witten, ``(2+1)-Dimensional Gravity as an Exactly Soluble System,'' \emph{Nucl. Phys. B} \textbf{311}, 46 (1988) doi:10.1016/0550-3213(88)90143-5

\bibitem{witten2007} E. Witten, ``Three-Dimensional Gravity Revisited,'' [arXiv:0706.3359 [hep-th]].

\end{thebibliography}

% =============================================================================
% 汇报（s05-wrapup.tex）
% * 覆盖：原文 p.32「A unified view.」至 p.36 §5 末；参考文献 [1]--[32]（p.36--38）
% * 页锚点：p.32--p.38 共 7 个，p.33/p.34/p.36 三处原文断页在句中，锚点就近放在分句处
% * 编号公式：4 个，(5.1)--(5.4)，\tag 逐字照原文；(5.4) 的波浪箭头按 amssymb 用 \leadsto
%   （lint.py 的 \left/\right 计数会把 \rightsquigarrow 误计为 \right，故未用后者）
% * 译注：8 条（易错点 4、含糊点 3、物理图像 1）
% * 文献：32 条，键名与顺序照 SPEC §8；卷号/页码/doi/arXiv 号逐条对 p.36--38 图像核过，
%   未见缺漏或数字错字。原文自身笔误照录并在下方登记。
% * 原文笔误：[23] 期刊缩写作「Fuct. Anal. Appl.」（应为 Funct.）；[22] 作者名作
%   「N. Woodhouse」（通行作 N. M. J. Woodhouse）；[9] 作者名后为句点而非逗号。
% * 记号：群名按原文排成直立体 \mathrm{SU}/\mathrm{ISO}/\mathrm{SO}，与 s04a/s04b 一致
%   （s01 用斜体，Lead 可统一）；(5.3) 的配分函数按 p.34 图像用 \mathcal{Z}，与 s04b 一致。
% * 存疑：q=e^{2πi/k}（树图）与 q=exp(2πi/(k+N))（正确）作者先给后改，已照原文层次
%   译出；c=kd/(k+N) 中的 k 已是重整化能级，与树图 k 同符号，已在译注点明。
% =============================================================================
